% Part: lambda-calculus
% Chapter: introduction
% Section: primitive-recursion

\documentclass[../../../include/open-logic-section]{subfiles}

\begin{document}

\olfileid{lam}{int}{pr}
\olsection{\usetoken{S}{lambda definable} অপেক্ষকগুলি আদিম পুনরাবৃত্তির অধীনে বদ্ধ}

আদিম পুনরাবৃত্তির ক্ষেত্রে এসে অবশেষে আমাদের কিছুটা কাজ করতে হবে। ধাপে
ধাপে এগোতে হবে। আগের মতোই ধরি, আমাদের কাছে ইতিমধ্যে এমন $G'$ ও~$H'$
পদ আছে, যা যথাক্রমে $g$ ও~$h$ অপেক্ষককে !!{lambda define} করে; আমরা এমন
একটি $F$ পদ চাই, যা নিচের সমীকরণদুটি দিয়ে সংজ্ঞায়িত অপেক্ষক~$f$-কে
!!{lambda define}s:
\begin{align*}
f(0, \vec z) & = g(\vec z) \\
f(x+1, \vec z) & = h(x, f(x,\vec z), \vec z).
\end{align*}
তাই সাধারণভাবে, $G'$ ও~$H'$ ল্যাম্বডা পদ দুটি দেওয়া থাকলে এমন একটি
পদ~$F$ খুঁজে পেলেই যথেষ্ট, যাতে
\begin{align*}
F(\num{0}, \vec z) & \equiv G'(\vec z) \\
F(\overline{n+1}, \vec z) & \equiv H'(\num{n}, F(\num{n}, \vec z), \vec z)
\end{align*}
প্রতিটি স্বাভাবিক সংখ্যা~$n$-এর জন্য সত্য হয়। $G'$ ও~$H'$ যথাক্রমে
$g$ ও~$h$-কে !!{lambda define} করে বলে, $\vec z$-এর জায়গায়
$\num{\vec m}$ সংখ্যাপদগুলি বসালে $F(\num{n+1}, \num{\vec m})$
সঠিক উত্তরের স্বাভাবিক রূপে হ্রাস পাবে।

কিন্তু এর জন্য এমন একটি পদ~$F$ খুঁজে পেলেই যথেষ্ট, যা
\begin{align*}
F(\num 0) & \equiv G \\
F(\overline {n+1}) & \equiv H(\num{n},F(\num{n}))
\intertext{প্রতিটি স্বাভাবিক সংখ্যা $n$-এর জন্য পূরণ করে, যেখানে}
G & = \lambd[\vec z][G'(\vec z)] \text{ এবং}\\
H(u,v) & = \lambd[\vec z][H'(u,v(\vec z),\vec z)].
\end{align*}
% BN-SRC-271 TARGET-CORRECTION-BEGIN
\par\smallskip\noindent\textnormal{\small\textbf{উৎস-সংশোধন (BN-SRC-271):} হিমায়িত অনুচ্ছেদে মূল সংজ্ঞাকারী পদদুটিকে প্রথমে~$G,H$ এবং পরক্ষণেই~$G',H'$ বলা হয়েছে; চাওয়া ফল-পদটিকেও ভুল করে~$H$ বলা হয়েছে; তদুপরি প্রথম প্রদর্শনে নিচে সংজ্ঞায়িত আবৃত~$G,H$ আগাম ব্যবহার করা হয়েছে। বাংলা পাঠে মূল পদদুটি সর্বত্র~$G',H'$, চাওয়া পদ~$F$, এবং~$G,H$ কেবল নিচের আবৃত পদদুটি।} % BN-SRC-271 TARGET-CORRECTION-END
% BN-SRC-272 TARGET-CORRECTION-BEGIN
\par\smallskip\noindent\textnormal{\small\textbf{উৎস-সংশোধন (BN-SRC-272):} হিমায়িত প্রথম পুনরাবৃত্তি-সমীকরণে~$h$-এর প্রথম আর্গুমেন্টে পুনরাবৃত্তি-সূচক~$x$-এর বদলে অসংগত~$z$ আছে। বাংলা সূত্রে~$h(x,f(x,\vec z),\vec z)$ পুনরুদ্ধার করা হয়েছে।} % BN-SRC-272 TARGET-CORRECTION-END
% BN-SRC-273 TARGET-CORRECTION-BEGIN
\par\smallskip\noindent\textnormal{\small\textbf{উৎস-সংশোধন (BN-SRC-273):} আবৃত ধাপ-পদে~$v$ হলো~$F(\num n)$, অর্থাৎ অবশিষ্ট~$\vec z$ আর্গুমেন্টের অপেক্ষক; হিমায়িত~$v(u,\vec z)$ তাই পুনরাবৃত্তি-সূচকটি দ্বিতীয়বার প্রয়োগ করে। বাংলা সূত্রে~$v(\vec z)$ লেখা হয়েছে, ফলে তা আগের~$F(\num n,\vec z)$ দফার ঠিক সমতুল্য।} % BN-SRC-273 TARGET-CORRECTION-END
অন্যভাবে বললে, ল্যাম্বডা-কৌশল ব্যবহার করে অতিরিক্ত পরামিতি $\vec z$ নিয়ে
আলাদা করে চিন্তা এড়ানো যায়—সেগুলি ল্যাম্বডা সংকেতলিপির মধ্যেই শোষিত হয়।

$F$ পদটি সংজ্ঞায়িত করার আগে ক্রমযুগল নিয়ে কাজ করার একটি ব্যবস্থা
দরকার। পরের লেমাটি সেই ব্যবস্থা দেয়।

\begin{lem}
এমন একটি ল্যাম্বডা পদ~$D$ আছে, যাতে যেকোনো দুটি ল্যাম্বডা পদ $M$ ও~$N$-এর
জন্য $D(M,N)(\num{0}) \red M$ এবং $D(M,N)(\num{1}) \red N$।
\end{lem}

\begin{proof}
প্রথমে ল্যাম্বডা পদ~$K$-কে সংজ্ঞায়িত করো:
\[
K(y) = \lambd[x][y].
\]
অর্থাৎ $K$ হলো $\lambd[y][\lambd[x][y]]$ পদটি। অন্যভাবে দেখলে, প্রতিটি
$M$-এর জন্য $K(M)$ এমন একটি ধ্রুবক অপেক্ষক, যা যেকোনো নিবেশে~$M$ ফেরত
দেয়।

এবার $D(x,y,z)$-কে $D(x,y,z) = z (K(y))x$ দিয়ে সংজ্ঞায়িত করো। তাহলে
\begin{align*}
D(M,N,\num 0) & \red \num 0 (K(N)) M \red M \text{ এবং}\\
D(M,N,\num 1) & \red \num 1 (K(N)) M \red K(N) M \red N,
\end{align*}
যেমনটি দরকার ছিল।
\end{proof}

মূল ধারণা হলো, $D(M,N)$ ক্রমযুগল $\tuple{M, N}$-কে উপস্থাপন করে; আর
$P$-কে এমন একটি ক্রমযুগলের উপস্থাপক ধরা হলে $P(\num 0)$ ও $P(\num 1)$
যথাক্রমে বাঁ ও ডান অভিক্ষেপ $(P)_0$ ও~$(P)_1$-কে উপস্থাপন করে। নিচে আমরা
শেষের সংকেতলিপিদুটি ব্যবহার করব।

\begin{lem}
!!{lambda definable} অপেক্ষকগুলি আদিম পুনরাবৃত্তির অধীনে বদ্ধ।
\end{lem}

\begin{proof}
আমাদের দেখাতে হবে, যেকোনো $G$ ও~$H$ পদ দেওয়া থাকলে এমন একটি পদ~$F$
পাওয়া যায়, যাতে
\begin{align*}
F(\num{0}) & \equiv G \\
F(\num{n+1}) & \equiv H(\num{n}, F(\num{n}))
\end{align*}
প্রতিটি স্বাভাবিক সংখ্যা~$n$-এর জন্য সত্য হয়। মোটামুটি ধারণাটি হলো,
সংখ্যাপদগুলিকে পুনরাবর্তক হিসেবে ব্যবহার করে নিচের \emph{ক্রমযুগলগুলির}
অনুক্রম গণনা করা:
\[
\tuple{\num 0, F(\num 0)}, \tuple{\num 1, F(\num 1)}, \dots,
\]
লক্ষ করো, প্রথম ক্রমযুগলটি ঠিক $\tuple{\num 0, G}$। কোনো ক্রমযুগল
$\tuple{\num{n}, F(\num{n})}$ দেওয়া থাকলে পরের ক্রমযুগল
$\tuple{\num{n+1}, F(\num{n+1})}$-এর সমতুল্য হওয়ার কথা
$\tuple{\num{n+1}, H(\num{n}, F(\num{n}))}$-এর। এমন এক-ধাপের রূপান্তর
ঘটায় এমন একটি ল্যাম্বডা পদ~$T$ আমরা তৈরি করব।

বিস্তারিত হলো এই। $T(u)$-কে সংজ্ঞায়িত করো:
\[
T(u) = \tuple{S((u)_0), H((u)_0,(u)_1)}.
\]
এখন সহজেই যাচাই করা যায়, যেকোনো সংখ্যা~$n$-এর জন্য
\[
T(\tuple{\num{n}, M}) \red \tuple{\num{n+1}, H(\num{n}, M)}.
\]
উপরের ইঙ্গিতমতো, $G$ ও~$H$ দেওয়া থাকলে $F(u)$-কে সংজ্ঞায়িত করো:
\[
F(u) = (u(T,\tuple{\num 0, G}))_1.
\]
অর্থাৎ নিবেশ~$\num{n}$-এ $F$, $\tuple{\num 0, G}$-এর উপর $T$-কে
$n$ বার পুনরাবৃত্ত করে, তারপর দ্বিতীয় উপাংশটি ফেরত দেয়। শুরুতে পাই
\begin{enumerate}
\item $\num{0} (T, \tuple{\num 0, G}) \equiv \tuple{\num{0}, G}$
\item $F(\num{0}) \equiv G$
\end{enumerate}
$n$-এর উপর আরোহ প্রয়োগ করে দেখানো যায় যে প্রতিটি স্বাভাবিক সংখ্যার
জন্য নিচের কথাদুটি সত্য:
\begin{enumerate}
\item $\num{n+1}(T, \tuple{\num{0}, G}) \equiv \tuple{\num{n+1}, F(\num{n+1})}$
\item $F(\num{n+1}) \equiv H(\num{n}, F(\num{n}))$
\end{enumerate}
দ্বিতীয় দফার জন্য পাই
\begin{align*}
F(\num{n+1}) & \red (\num{n+1}(T, \tuple{\num 0, G}))_1 \\
& \equiv (T(\num{n} (T, \tuple{\num 0, G})))_1 \\
& \equiv (T(\tuple{\num{n}, F(\num{n})}))_1 \\
& \equiv (\tuple{\num{n+1}, H(\num{n}, F(\num{n}))})_1 \\
& \equiv H(\num{n}, F(\num{n})).
\end{align*}
এখানে শেষের আগের পঙ্‌ক্তিতে আরোহানুমান ব্যবহার করেছি। প্রথম দফার জন্য
পাই
\begin{align*}
\num{n+1} (T, \tuple{\num 0, G}) &
\equiv T(\num{n} (T, \tuple{\num 0, G})) \\
& \equiv T( \tuple{\num{n}, F(\num{n})}) \\
& \equiv \tuple{\num{n+1}, H(\num{n}, F(\num{n}))} \\
& \equiv \tuple{\num{n+1}, F(\num{n+1})}.
\end{align*}
এখানে শেষ পঙ্‌ক্তিতে দ্বিতীয় দফাটি ব্যবহার করেছি। অতএব আমরা দেখিয়েছি
$F(\num 0) \equiv G$ এবং প্রতিটি $n$-এর জন্য $F(\num {n+1}) \equiv H(\num
n, F(\num{n}))$; ঠিক এটিই আমাদের দরকার ছিল।
\end{proof}

\end{document}
